% Source-aligned English for printed p.142.
\SourcePage{142}

\LayerTarget{tohoku:thm:2.2.2:p2}
\noindent\textit{is necessary and sufficient that $F$ satisfy the following
conditions: $F$ is half-exact, and for $P\subset Q\subset R$ in $\mathbf{C}$, the
kernel of $F(Q/P)\to F(R/P)$ is contained in the image of
$F(Q)\to F(Q/P)$. These conditions are always satisfied if $F$
is left exact or right exact.}\textsuperscript{2bis)}\SourceFootnote{2 bis}{Added in proof. This condition is also
automatically satisfied if every object of $C$ is isomorphic to a subobject
of an injective object; see [6, Chap. III].}

\LayerTarget{tohoku:thm:2.2.2:proof:p1}
The proof is essentially contained in [6, Chap. III]. For the first part it
suffices to prove the existence of $S^1F$. Let $A\in \mathbf{C}$ and consider an exact
sequence
$0\to A\to M\to A'\to0$ in
which the first morphism is an injective effacement of $A$. Set
$S^1F(A)=F(A')/\mathrm{Im}(F(M))$. As in [6], the latter object is independent
of the particular exact sequence chosen, up to canonical isomorphism, and can
be regarded as a functor of $A$. The definition of the boundary
homomorphism, verification of axioms (i) and (ii) of 2.1, and verification
that the resulting $\partial$-functor $(F,S^1F)$ is universal are likewise
standard. We shall also omit the proof of the exactness criterion.
\LayerTarget{tohoku:thm:2.2.2:dual}
--- We record the dual statement. \textit{If every object of $\mathbf{C}$ has a
projective effacement, then the satellites $S^iF$ for $i\leqq0$ exist and are
co-effaceable functors for $i<0$; the condition for the $\partial$-functor
$(S^iF)_{i\leqq0}$ to be exact is the same as in Theorem \textup{2.2.2}.} Consequently,
if every object of $\mathbf{C}$ has both an injective effacement and a projective
effacement, every covariant additive functor $F$ has satellites $S^iF$ for
all $i$, and the $\partial$-functor $(S^iF)$ is exact if and only if $F$
satisfies the condition in Theorem 2.2.2. If $F$ is contravariant, the
cases $i\leqq0$ and $i\geqq0$ in the foregoing statements must be interchanged,
and the exactness condition replaced by its dual.

\LayerTarget{tohoku:remark:satellites-small-skeleton:p1}
\indent\textsc{Remark.} We mention another case, very different from that
of Theorem 2.2.2, in which one can construct the satellite functors of an
arbitrary functor from $\mathbf{C}$ to $\mathbf{C}'$. Suppose that there is a \textit{set}
$\mathbf{C}_0\subset \mathbf{C}$ such that every $A\in \mathbf{C}$ is isomorphic to an object of $\mathbf{C}_0$,
and suppose that $\mathbf{C}'$ is an abelian category in which infinite direct sums
exist. Then the satellites $S^iF$, $i\geqq0$, exist for every additive functor
$F:\mathbf{C}\to \mathbf{C}'$. Moreover, if $\mathbf{C}'$ satisfies axiom AB 5), see 1.5,
and $F$ satisfies the condition at the end of Theorem 2.2.2, then the
$\partial$-functor $(S^iF)_{i\geqq0}$ is \textit{exact}. The category of abelian groups,
in particular, satisfies AB 5), so the preceding result can be applied to the
functor $\mathrm{Hom}(A,B)$ with values in abelian groups. One can thereby
define $\mathrm{Ext}^i(A,B)$ as its satellites, regarding
$\mathrm{Hom}(A,B)$ either as a covariant functor of $B$ or as a contravariant
functor of $A$. It would still remain to prove that the two procedures give
the same result. The condition imposed on $\mathbf{C}$ holds for categories whose
objects satisfy certain finiteness conditions, and in which infinite direct
sums therefore do not generally exist. For example, it holds for the
abelian category of not necessarily connected algebraic groups defined over
a fixed field $k$ of characteristic zero which are complete as algebraic
varieties and abelian as groups, that is, the category of algebraic groups

\TohokuContinuation{143}
